% Storyline 3 (final), paragraph 5: findings and broad implication.
We find that the direction of revision, rather than the amount of evidence it requires, decides whether private observers agree with one another and still tell cooperators from defectors. The beneficial direction mirrors the kind of misjudgement that private assessment produces under each norm, and rules selected solely by these reputational criteria make cooperation stable and easy to reach under four classical norms. The same amount of inertia placed in the wrong direction can eliminate cooperation altogether. Rules that protect cooperators from disputed evidence, however, can also shelter opportunists who defect only occasionally, a weakness shared by the aggregation and score-based rules we test. Letting suspicion carry partial consequences restores deterrence while cooperation remains high. Sustaining cooperation through private reputations is therefore not only a matter of what individuals know about each other, but also of how they revise what they believe.
